% GENERATED FILE. Edit canonical lessons, then run npm run generate:books.
% canonical-source-hash: fnv1a64:c3524ae8febc6239
% canonical-lessons: MR-C546-sankalpa, MR-C546-nirdhar, MR-C546-irada, MR-C546-sapath-ghene, MR-C546-dhajne

\chapter{Firm Intentions}
\label{ch:mr-firm-intentions}
\hlchaptermodality{546}

\begin{chapteropening}
\textbf{By the end of this chapter:} \emph{I can talk about resolutions, oaths and daring.}

Complete the last lesson of chapter 546: I can talk about resolutions, oaths and daring.
\end{chapteropening}

\section[\texorpdfstring{\mr{संकल्प} (saṅkalpa)}{संकल्प (saṅkalpa)}]{\mr{संकल्प} (saṅkalpa) — a resolution; a vow}

\label{lesson:MR-C546-sankalpa}



\begin{quote}
Before the new one: say the Marathi for to conjecture, then the Marathi for to forecast.
\end{quote}

\subsection*{You'll want to know: \mr{संकल्प}}
\textbf{\mr{संकल्प}} — \emph{saṅkalpa} — \textquotedblleft{}a resolution; a vow\textquotedblright{}.

\begin{quote}
\mr{नवीन} \mr{वर्षी} \mr{मी} \mr{रोज} \mr{चालण्याचा} \mr{संकल्प} \mr{केला}.
\end{quote}

\begin{quote}
\emph{navīn varṣī mī roj cālṇyācā saṅkalpa kelā.}
\end{quote}

\begin{quote}
In the new year I resolved to walk every day.
\end{quote}

Masculine.

\mr{संकल्प} \mr{करणे}, to make a resolution.

\begin{cousinweb}
\textbf{\mr{संकल्प}} (\emph{saṅkalpa}) is a tatsama, taken from Sanskrit \emph{saṁkalpa}, \textquotedblleft{}resolve, intention\textquotedblright{} (Monier-Williams).
\end{cousinweb}

\begin{grammarlens}[title={What you've built}]
One more word for saying what happens if.
\end{grammarlens}

\subsection*{Your turn}
\begin{itemize}
\raggedright
  \item \emph{Say it:} \emph{saṅkalpa}
  \item \emph{Say it:} \emph{saṅkalpa}, once more
  \item \emph{Say it:} \emph{navīn varṣī mī roj cālṇyācā saṅkalpa kelā.} — In the new year I resolved to walk every day.
  \item \emph{Recall:} say the Marathi for to conjecture, then the Marathi for to forecast, then say \emph{saṅkalpa} again
\end{itemize}

\begin{tcolorbox}[breakable,colback=teal!4,colframe=teal!35!black,title={Before you move on}]
What does \mr{संकल्प} mean? (\textquotedblleft{}A resolution; a vow\textquotedblright{}.) Say it once more.
\end{tcolorbox}
\section[\texorpdfstring{\mr{निर्धार} (nirdhār)}{निर्धार (nirdhār)}]{\mr{निर्धार} (nirdhār) — determination}

\label{lesson:MR-C546-nirdhar}



\begin{quote}
Before the new one: say the Marathi for to forecast, then the Marathi for a resolution.
\end{quote}

\subsection*{You'll want to know: \mr{निर्धार}}
\textbf{\mr{निर्धार}} — \emph{nirdhār} — \textquotedblleft{}determination\textquotedblright{}.

\begin{quote}
\mr{परीक्षा} \mr{पास} \mr{होण्याचा} \mr{तिचा} \mr{निर्धार} \mr{पक्का} \mr{होता}.
\end{quote}

\begin{quote}
\emph{parīkṣā pās hoṇyācā ticā nirdhār pakkā hotā.}
\end{quote}

\begin{quote}
Her determination to pass the exam was firm.
\end{quote}

Masculine.

\begin{cousinweb}
\textbf{\mr{निर्धार}} (\emph{nirdhār}) is a tatsama, taken from Sanskrit \emph{nirdhāra}, \textquotedblleft{}settling, certainty\textquotedblright{} (Monier-Williams).
\end{cousinweb}

\begin{grammarlens}[title={What you've built}]
One more word for saying what happens if.
\end{grammarlens}

\subsection*{Your turn}
\begin{itemize}
\raggedright
  \item \emph{Say it:} \emph{nirdhār}
  \item \emph{Say it:} \emph{nirdhār}, once more
  \item \emph{Say it:} \emph{parīkṣā pās hoṇyācā ticā nirdhār pakkā hotā.} — Her determination to pass the exam was firm.
  \item \emph{Recall:} say the Marathi for to forecast, then the Marathi for a resolution, then say \emph{nirdhār} again
\end{itemize}

\begin{tcolorbox}[breakable,colback=teal!4,colframe=teal!35!black,title={Before you move on}]
What does \mr{निर्धार} mean? (\textquotedblleft{}Determination\textquotedblright{}.) Say it once more.
\end{tcolorbox}
\section[\texorpdfstring{\mr{इरादा} (irādā)}{इरादा (irādā)}]{\mr{इरादा} (irādā) — an intention}

\label{lesson:MR-C546-irada}



\begin{quote}
Before the new one: say the Marathi for a resolution, then the Marathi for determination.
\end{quote}

\subsection*{You'll want to know: \mr{इरादा}}
\textbf{\mr{इरादा}} — \emph{irādā} — \textquotedblleft{}an intention\textquotedblright{}.

\begin{quote}
\mr{पुढच्या} \mr{वर्षी} \mr{परदेशी} \mr{जाण्याचा} \mr{माझा} \mr{इरादा} \mr{आहे}.
\end{quote}

\begin{quote}
\emph{puḍhcyā varṣī pardeśī jāṇyācā mājhā irādā āhe.}
\end{quote}

\begin{quote}
I intend to go abroad next year.
\end{quote}

Masculine.

\begin{cousinweb}
\textbf{\mr{इरादा}} (\emph{irādā}) is a loan through Persian usually traced to Arabic \emph{irāda}, \textquotedblleft{}will, wish\textquotedblright{}.
\end{cousinweb}

\begin{grammarlens}[title={What you've built}]
One more word for saying what happens if.
\end{grammarlens}

\subsection*{Your turn}
\begin{itemize}
\raggedright
  \item \emph{Say it:} \emph{irādā}
  \item \emph{Say it:} \emph{irādā}, once more
  \item \emph{Say it:} \emph{puḍhcyā varṣī pardeśī jāṇyācā mājhā irādā āhe.} — I intend to go abroad next year.
  \item \emph{Recall:} say the Marathi for a resolution, then the Marathi for determination, then say \emph{irādā} again
\end{itemize}

\begin{tcolorbox}[breakable,colback=teal!4,colframe=teal!35!black,title={Before you move on}]
What does \mr{इरादा} mean? (\textquotedblleft{}An intention\textquotedblright{}.) Say it once more.
\end{tcolorbox}
\section[\texorpdfstring{\mr{शपथ} \mr{घेणे} (śapath gheṇe)}{शपथ घेणे (śapath gheṇe)}]{\mr{शपथ} \mr{घेणे} (śapath gheṇe) — to take an oath}

\label{lesson:MR-C546-sapath-ghene}



\begin{quote}
Before the new one: say the Marathi for determination, then the Marathi for an intention.
\end{quote}

\subsection*{You'll want to know: \mr{शपथ} \mr{घेणे}}
\textbf{\mr{शपथ} \mr{घेणे}} — \emph{śapath gheṇe} — \textquotedblleft{}to take an oath\textquotedblright{}.

\begin{quote}
\mr{नव्या} \mr{महापौरांनी} \mr{सभेत} \mr{शपथ} \mr{घेतली}.
\end{quote}

\begin{quote}
\emph{navyā mahāpaurānnī sabhet śapath ghetlī.}
\end{quote}

\begin{quote}
The new mayor took the oath at the meeting.
\end{quote}

Past, with -\mr{ने} (he …): \mr{त्याने} \mr{शपथ} \mr{घेतली}, \emph{tyāne śapath ghetlī}.

\begin{cousinweb}
\textbf{\mr{शपथ} \mr{घेणे}} (\emph{śapath gheṇe}) is the tatsama \mr{शपथ}, from Sanskrit \emph{śapatha}, \textquotedblleft{}a curse, an oath\textquotedblright{} (Monier-Williams), with \mr{घेणे}.
\end{cousinweb}

\begin{grammarlens}[title={What you've built}]
One more word for saying what happens if.
\end{grammarlens}

\subsection*{Your turn}
\begin{itemize}
\raggedright
  \item \emph{Say it:} \emph{śapath gheṇe}
  \item \emph{Say it:} \emph{śapath gheṇe}, once more
  \item \emph{Say it:} \emph{navyā mahāpaurānnī sabhet śapath ghetlī.} — The new mayor took the oath at the meeting.
  \item \emph{Recall:} say the Marathi for determination, then the Marathi for an intention, then say \emph{śapath gheṇe} again
\end{itemize}

\begin{tcolorbox}[breakable,colback=teal!4,colframe=teal!35!black,title={Before you move on}]
What does \mr{शपथ} \mr{घेणे} mean? (\textquotedblleft{}To take an oath\textquotedblright{}.) Say it once more.
\end{tcolorbox}
\section[\texorpdfstring{\mr{धजणे} (dhajṇe)}{धजणे (dhajṇe)}]{\mr{धजणे} (dhajṇe) — to dare}

\label{lesson:MR-C546-dhajne}



\begin{quote}
Before the new one: say the Marathi for an intention, then the Marathi for to take an oath.
\end{quote}

\subsection*{You'll want to know: \mr{धजणे}}
\textbf{\mr{धजणे}} — \emph{dhajṇe} — \textquotedblleft{}to dare\textquotedblright{}.

\begin{quote}
\mr{अंधारात} \mr{एकट्याने} \mr{जायला} \mr{तो} \mr{धजला} \mr{नाही}.
\end{quote}

\begin{quote}
\emph{andhārāt ekṭyāne jāylā to dhajlā nāhī.}
\end{quote}

\begin{quote}
He did not dare to go alone in the dark.
\end{quote}

Past (he …): \mr{तो} \mr{धजला}, \emph{to dhajlā}.

It takes the -\mr{ायला} form: \mr{बोलायला} \mr{धजणे}, to dare to speak.

\begin{cousinweb}
\textbf{\mr{धजणे}} (\emph{dhajṇe}) is a Marathi verb heard mostly with a negative; its origin is not settled.
\end{cousinweb}

\begin{grammarlens}[title={What you've built}]
One more word for saying what happens if.
\end{grammarlens}

\subsection*{Your turn}
\begin{itemize}
\raggedright
  \item \emph{Say it:} \emph{dhajṇe}
  \item \emph{Say it:} \emph{dhajṇe}, once more
  \item \emph{Say it:} \emph{andhārāt ekṭyāne jāylā to dhajlā nāhī.} — He did not dare to go alone in the dark.
  \item \emph{Recall:} say the Marathi for an intention, then the Marathi for to take an oath, then say \emph{dhajṇe} again
\end{itemize}

\begin{tcolorbox}[breakable,colback=teal!4,colframe=teal!35!black,title={Before you move on}]
What does \mr{धजणे} mean? (\textquotedblleft{}To dare\textquotedblright{}.) Say it once more.
\end{tcolorbox}
